\begin{textsample}{POS Dim 3 – rn – Score 19.00 – rn/rn\_inf\_6.txt}  \label{ex:f3_pos_098}
2.
A PRIMEIRA \textbf{INFÂNCIA}, \textbf{CRIANÇAS} E CULTURAS \textbf{INFANTIS} [... ] as \textbf{crianças} precisam ser concebidas como pessoas concretas – e não seres ideais e abstratos ; como pessoas contemporâneas – já existentes em seu aqui e agora e não como \textbf{seres} “ do futuro ” ; com direitos iguais aos de qualquer indivíduo humano ; com necessidades e potencialidades humanas – dependentes dos outros e capazes de interagir, de aprender, de se desenvolver e de produzir cultura – de \textbf{inventar}, criar outros modos de se relacionar com o mundo e de entendê -lo ; como seres integrais que se manifestam em sua completude e não de modo segmentado, em partes separadas ; que, por serem humanas, as \textbf{crianças} são semelhantes em todas as culturas, mas, por serem \textbf{seres} da cultura, são singulares ; cada uma é única.
São muitas e diversas as \textbf{crianças}. ( UBARANA ; LOPES, 2012, pp. 2-3 ) São as \textbf{crianças} diversas e singulares, concretas e contemporâneas que habitam as \textbf{instituições} de Educação \textbf{Infantil} na relação com os profissionais, as famílias, os espaços e as propostas e práticas de considerá -las e incluí -las nos processos de planejamento e avaliação.
Desse modo, objetiva -se, neste capítulo, refletir sobre as \textbf{infâncias}, as \textbf{crianças} e suas especificidades e culturas.
Pois compreende -se que toda proposta precisa contextualizar os sujeitos e suas práticas culturais na definição de objetivos e experiências do currículo.
CONCEPÇÕES DE \textbf{INFÂNCIAS} E \textbf{CRIANÇAS} A \textbf{infância} é compreendida como construção histórica, social e cultural, portanto, suscetível a mudanças sempre que grandes transformações culturais ocorrem historicamente.
Desse modo, considera -se a \textbf{infância} como tempo e condição de vida das \textbf{crianças}, como categoria geracional das sociedades ( geração de indivíduos entre zero a 10-12 anos, com especificidades relativas à primeira e segunda \textbf{infância} ) – constituída por dimensões : biológica-etária ( relativa às características humanas ) ; sócio-histórico-cultural ( relativa às condições concretas de vida das \textbf{crianças} ), bem como aspectos que marcam as diferentes \textbf{infâncias}, como a nacionalidade, a classe social, o gênero, a etnia, a religião, entre outros. [... ] a \textbf{infância} é, ao mesmo tempo, uma “ categoria geracional ” – uma classe à qual pertencem os humanos de uma certa idade –, mas também um grupo social composto por sujeitos com características próprias – as \textbf{crianças} – que vivem e agem no mundo e deixam suas \textbf{marcas}, produzem modos próprios de ser e estar no mundo, de entender, de agir, de pensar, de sentir, de dizer.
Esses modos são marcados pela ludicidade, pela imaginação, pela \textbf{brincadeira} – que se distinguem dos modos “ \textbf{adultos} ” de verem o mundo, as \textbf{crianças} e a si próprios. ( UBARANA ; LOPES, 2012, p. 3 ) Sendo assim, as \textbf{infâncias} são semelhantes nesses modos de ser e estar no mundo, mas são diversas de acordo com condições que marcam a vida das \textbf{crianças}.
Toda \textbf{criança} vive sua própria \textbf{infância} – segundo suas próprias condições de vida.
É preciso, portanto, levar em conta o tempo, as necessidades e os modos das \textbf{crianças} viverem suas \textbf{infâncias}. “ É necessário conhecer as representações de \textbf{infância}, compreendendo as \textbf{crianças} reais, concretas, nas relações sociais, como produtoras de cultura. ” ( KUHLMANN JR., 2015, p. 30 ).
Em um estado como o nosso, diverso nos seus modos de organização e nas particularidades culturais de cada localidade, são diversas as \textbf{infâncias} vividas pelas \textbf{crianças}.
Algumas \textbf{crianças} vivem em contextos urbanos e outras, no campo ou no litoral.
Os modos de \textbf{brincar} variam de acordo com as práticas culturais e as condições de vida : algumas \textbf{crianças} transformam em \textbf{brinquedos} os elementos da natureza ou de sua própria casa e outras fazem uso, como \textbf{brinquedo}, dos recursos tecnológicos ; umas têm liberdade e um vasto espaço para correr e explorar e outras precisam limitar -se à sua própria casa ; umas têm tempo definido na \textbf{rotina} para suas escolhas e outras usam do seu tempo para estudar e trabalhar, sem ter direito a escolher o que fazer.
Algumas \textbf{crianças} têm acesso à Educação \textbf{Infantil} desde \textbf{bebês}, outras ingressam aos 2 ou 3 anos de idade.
As práticas cotidianas vivenciadas nas \textbf{instituições} de Educação \textbf{Infantil} também são, portanto, constitutivas da \textbf{infância} de cada \textbf{criança}.
O Currículo vivido nas \textbf{instituições} é considerado como fundante das \textbf{infâncias} das \textbf{crianças} – das condições de interações e mediações às quais elas têm acesso para aprenderem e se desenvolverem integralmente.
APROFUNDANDO O TEMA Sociologia da \textbf{Infância} Apresenta -se “ um novo entendimento da \textbf{infância} e das \textbf{crianças} ”, considerando que a \textbf{infância} é uma construção social, elaborada para e pelas \textbf{crianças}, em um conjunto ativamente negociado de relações sociais.
Embora a \textbf{infância} seja um fato biológico, a maneira como ela é entendida é determinada socialmente [... ] é sempre contextualizada em relação ao local e à cultura, variando segundo a classe, o gênero e outras condições socioeconômicas.
Por isso, não há uma \textbf{infância} natural nem universal, e nem uma \textbf{criança} natural ou universal, mas muitas \textbf{infâncias} e \textbf{crianças}. ( DAHLBERG ; MOSS ; PENCE, 2003, p. 71 ).
Esta construção se fundamenta especialmente no discurso do “ novo paradigma da sociologia da \textbf{infância} ” que, segundo Sarmento ( 2008 ), tem origem já na década de 1990 com os “ aspectos-chave do paradigma ” definidos por Prout e James sobre a investigação sociológica da \textbf{infância}, dos quais cita -se aqui : 1.
A \textbf{infância} é entendida como construção social.
Como tal, isso indica um quadro interpretativo para a contextualização dos primeiros anos da vida humana.
A \textbf{infância}, sendo distinta da imaturidade biológica, não é uma forma natural nem universal dos grupos humanos, mas aparece como uma componente estrutural e cultural específica de muitas sociedades. 2.
A \textbf{infância} é uma variável de análise social.
Ela não pode nunca ser inteiramente divorciada de outras variáveis como a classe social, o gênero ou a pertença étnica.
A análise comparativa e multicultural revela uma variedade de \textbf{infâncias}, mais do que um fenômeno singular e universal. 3.
As relações sociais estabelecidas pelas \textbf{crianças} e suas culturas devem ser estudadas por seu próprio direito, independentemente da perspectiva e dos conceitos dos \textbf{adultos}. 4.
As \textbf{crianças} são e devem ser vistas como atores na construção e determinação de suas próprias vidas sociais, das vidas dos que as rodeiam e das sociedades em que vivem.
As \textbf{crianças} não são os sujeitos passivos de estruturas e processos sociais [... ] ( PROUT ; JAMES, 1990, p. 8-9 apud SARMENTO, 2008, pp. 23-24 ).
Considera -se que é a partir da Sociologia da \textbf{Infância} que se consolida a ideia de que a \textbf{criança} é um ator social com direitos, impulsionando dessa forma o seu reconhecimento como cidadão ativo com lugar nas esferas social, política e científica.
Esse percurso abre renovadas possibilidades de considerar as \textbf{crianças}, bem como as relações e as práticas sociais desenvolvidas com elas, no sentido de contestar o entendimento das \textbf{crianças} como objetos passivos das políticas e das práticas \textbf{adultas}, cuja cidadania é vista como um potencial e um estatuto a ser alcançado no futuro ( CUNHA, 2016 ).
As \textbf{crianças} são vistas como cidadãos com direitos, \textbf{membros} de um grupo social, agentes de suas próprias vidas ( embora não agentes livres ) e como co-construtores do conhecimento, identidade e cultura. “ A \textbf{criança} como forte, competente, inteligente, um pedagogo poderoso, capaz de produzir teorias \textbf{interessantes} e desafiadoras, compreensões, perguntas – e desde o \textbf{nascimento}, não em uma idade avançada quando ficaram prontos ” ( MOSS, 2005, p. 242 ).
Nesse sentido, concebe -se neste documento, a \textbf{criança} como : - Sujeito histórico, social e cultural – cuja singularidade vai se construindo nas condições de vida concreta que lhes são possibilitadas – em sua história de vida social e pessoal ; - Pessoa que aprende e se desenvolve mediante condições de interação e mediação social e simbólica – mediante participação do(s) outro(s) e da(s) linguagem(s) ; - Ser concreto e contemporâneo ( do tempo presente ) : já é uma pessoa, cuja identidade/subjetividade está em processo de constituição inicial ; - Sujeito de direitos, dentre eles, à educação ( que envolve \textbf{cuidado} ) como condição para que possa desenvolver suas potencialidades humanas.
Uma \textbf{criança} com uma voz para ser \textbf{ouvida}, mas compreendendo que \textbf{ouvir} é um processo interpretativo e que as \textbf{crianças} podem se fazer \textbf{ouvir} de muitas formas ( conhecidamente expresso em As Cem Linguagens da \textbf{Infância}, de Malaguzzi ). ( MOSS, 2005 ).
Nesse sentido, a partir do interesse da \textbf{criança}, a consideração da sua participação implica que a sua voz seja integrada nos processos de tomada de decisão nos assuntos que lhe dizem respeito, de forma a ultrapassar a ideia apresentada por Qvortrup ( 1999, p. 9 ) de que “ (... ) os \textbf{adultos} afirmam que as \textbf{crianças} devem ser \textbf{ouvidas}, mas na maioria das vezes são tomadas decisões, que vão ter consequências nas suas vidas, sem que as mesmas sejam levadas em conta ”.
Considera -se aqui que as \textbf{crianças} são as melhores informantes sobre as questões que lhes dizem respeito.
E, portanto, são legitimadas como atores sociais, em condições de explicitarem sua compreensão sobre algo, pelo descortinamento de seus modos de ser e estar no mundo, não só pela fala, mas por múltiplas outras formas de comunicação que possuem. ( CUNHA, 2016 ).
Nessa compreensão, as \textbf{crianças} de cada contexto do estado do Rio Grande do Norte são pessoas singulares que se constituem nas relações, interações e práticas culturais, com seus modos próprios de aprender, investigar e se desenvolver.
Para tanto, precisam ser observadas, \textbf{ouvidas} e consideradas no planejamento e desenvolvimento de currículos e práticas pedagógicas em \textbf{instituições} educativas.
Desse modo, como produção de sua história de vida, cada \textbf{criança} tem suas singularidades, características que as constituem como pessoa e lhe conferem uma contemporaneidade enquanto sujeito histórico e social.
No entanto, é possível falar de especificidades que diferenciam as \textbf{crianças} \textbf{pequenas} em relação aos \textbf{seres} humanos de outros ciclos de vida, no que se refere aos modos de se relacionar com o mundo, de compreendê -lo e de agir nele/sobre ele, de viver, sentir, olhar, falar, aprender, simbolizar, movimentar -se e produzir culturas \textbf{infantis} vinculadas, principalmente, à ludicidade.

% matched lemmas: adulto, bebê, brincadeira, brincar, brinquedo, criança, cuidado, infantil, infância, instituição, interessante, inventar, marca, membro, nascimento, ouvir, pequeno, rotina, seres
\end{textsample}
